\scanpage{71}
C'est trivial. \deleted{Cas particuliers :}

\deleted{\uline{Proposition} Dans le cas \deleted{\unsure{d'un X}} I' ou II'', si $\unsure{\mathbb W|W}^a(\unsure{\mathfrak X|\mathcal X})$ existe et}

\uline{Lemme 2} Soient $\mathfrak X_i$ des ouverts de $\mathfrak X/\mathbb W(S)$, on suppose \deleted{\inserted{que $\underline{\mathbb W}^a(\mathfrak X)$ existe}} que \deleted{les $\mathfrak X_i$ dans $\mathfrak X$ satisfont les conditions du lemme 1 (en particulier} les $\mathbb W^a(\mathfrak X_i)$ existent \deleted{\unsure{et}} \deleted{\unsure{et que}} et que les $\mathfrak X_i\cap\mathfrak X_j$ dans $\mathfrak X_i$ satisfont les \deleted{\unsure{mêmes}} conditions \inserted{du lemme 1} [donc $\underline{\mathbb W}^a(\mathfrak X_i\cap\mathfrak X_j)$ existe] et est un ouvert dans $\underline{\mathbb W}^a(\mathfrak X_i)$], \inserted{et kif-kif pour les $\mathfrak X_i\cap\mathfrak X_j\cap\mathfrak X_{\unsure{h|k}}$ dans $\mathfrak X_i\cap\mathfrak X_j$} On suppose

\textcircled{$\beta$} Pour tt $T/S$, et tout morphisme $\underline{\mathbb W}(T)\xrightarrow f\mathfrak X$, on a $T=\bigcup_i T_i$ (les $T_i$ ouverts \unsure{de|dans} $T$) tel que
\[
 f(\mathbb W(T_i))\subset\mathfrak X_i
\]
Sous ces conditions, $\underline{\mathbb W}^a(\mathfrak X)$ existe, et est la réunion des $\underline{\mathbb W}^a(\mathfrak X_i)$, qui s'identifient à des ouverts de $\underline{\mathbb W}^a(\mathfrak X)$
% archive sheet number: 70
